% Options for packages loaded elsewhere
\PassOptionsToPackage{unicode}{hyperref}
\PassOptionsToPackage{hyphens}{url}
\PassOptionsToPackage{dvipsnames,svgnames,x11names}{xcolor}
\documentclass[
  10pt,
  fontset=fandol]{ctexart}
\usepackage{xcolor}
\usepackage[margin=16mm]{geometry}
\usepackage{amsmath,amssymb}
\setcounter{secnumdepth}{-\maxdimen} % remove section numbering
\usepackage{iftex}
\ifPDFTeX
  \usepackage[T1]{fontenc}
  \usepackage[utf8]{inputenc}
  \usepackage{textcomp} % provide euro and other symbols
\else % if luatex or xetex
  \usepackage{unicode-math} % this also loads fontspec
  \defaultfontfeatures{Scale=MatchLowercase}
  \defaultfontfeatures[\rmfamily]{Ligatures=TeX,Scale=1}
\fi
\usepackage{lmodern}
\ifPDFTeX\else
  % xetex/luatex font selection
\fi
% Use upquote if available, for straight quotes in verbatim environments
\IfFileExists{upquote.sty}{\usepackage{upquote}}{}
\IfFileExists{microtype.sty}{% use microtype if available
  \usepackage[]{microtype}
  \UseMicrotypeSet[protrusion]{basicmath} % disable protrusion for tt fonts
}{}
\makeatletter
\@ifundefined{KOMAClassName}{% if non-KOMA class
  \IfFileExists{parskip.sty}{%
    \usepackage{parskip}
  }{% else
    \setlength{\parindent}{0pt}
    \setlength{\parskip}{6pt plus 2pt minus 1pt}}
}{% if KOMA class
  \KOMAoptions{parskip=half}}
\makeatother
\setlength{\emergencystretch}{3em} % prevent overfull lines
\providecommand{\tightlist}{%
  \setlength{\itemsep}{0pt}\setlength{\parskip}{0pt}}
\usepackage{amsmath,amssymb,mathtools,needspace}
\xeCJKDeclareCharClass{CJK}{"2032,"0400->"04FF,"2160->"216B,"2460->"2473,"25A0,"25C7,"25CB,"25CF}
\IfFontExistsTF{Arial Unicode MS}{\xeCJKsetup{AutoFallBack=true}\setCJKfallbackfamilyfont{\CJKrmdefault}{Arial Unicode MS}}{}
\setcounter{secnumdepth}{0}
\setlength{\emergencystretch}{2em}
\allowdisplaybreaks
\usepackage{bookmark}
\IfFileExists{xurl.sty}{\usepackage{xurl}}{} % add URL line breaks if available
\urlstyle{same}
\hypersetup{
  pdftitle={Gauss 公式的余项},
  colorlinks=true,
  linkcolor={Maroon},
  filecolor={Maroon},
  citecolor={Blue},
  urlcolor={Blue},
  pdfcreator={LaTeX via pandoc}}

\title{Gauss 公式的余项}
\author{}
\date{}

\begin{document}
\maketitle

\subsubsection{4.4.3 Gauss 公式的余项}\label{gauss-ux516cux5f0fux7684ux4f59ux9879}

\textbf{定理 4.5}　对于 Gauss 公式（4.4.1），其余项

\[
R(x)=\int_a^b f(x)\,\mathrm{d}x-\sum_{k=0}^{n}A_kf(x_k)
=\frac{f^{(2n+2)}(\xi)}{(2n+2)!}\int_a^b\omega^2(x)\,\mathrm{d}x,
\tag{4.4.5}
\]

这里 \(\omega(x)=(x-x_0)(x-x_1)\cdots(x-x_n)\)。

\textbf{证明}　以 \(x_0,x_1,\cdots,x_n\) 为节点构造次数不大于 \(2n+1\) 的多项式 \(H(x)\)，使其满足条件

\[
H(x_i)=f(x_i),\quad H'(x_i)=f'(x_i)\qquad(i=0,1,\cdots,n).
\]

我们知道，这样的 \(H(x)\) 称为 Hermite 插值多项式（见式（2.6.3）、（2.6.6））。由于 Gauss 公式具有 \(2n+1\) 次代数精度，它对于 \(H(x)\) 能准确成立，即

\[
\int_a^b H(x)\,\mathrm{d}x
=\sum_{k=0}^{n}A_kH(x_k)=\sum_{k=0}^{n}A_kf(x_k),
\]

故有

\[
\begin{aligned}
R(x)&=\int_a^b f(x)\,\mathrm{d}x-\sum_{k=0}^{n}A_kf(x_k)
=\int_a^b f(x)\,\mathrm{d}x-\int_a^b H(x)\,\mathrm{d}x\\
&=\int_a^b\frac{f^{(2n+2)}(\eta)}{(2n+2)!}\omega^2(x)\,\mathrm{d}x.
\end{aligned}
\]

再考虑到该函数 \(\omega^2(x)\) 在 \([a,b]\) 上保号，应用积分中值定理得式（4.4.5）。证毕。

\end{document}
